\documentclass[10pt,a4paper]{amsart}
\usepackage[margin=19mm]{geometry}
\usepackage{amsmath,amssymb,mathtools}
\usepackage[scr=rsfs,cal=euler]{mathalfa}
\usepackage{xfrac}
\usepackage[unicode,hidelinks,bookmarks=false]{hyperref}
\newcommand{\ie}{i.e.\ }
\newcommand{\cf}{cf.\ }
\newcommand{\resp}{respectively\ }
\newcommand{\Subsection}{\subsection}
\providecommand{\nobreakdash}{\nobreak\hbox{-}}
\setlength{\emergencystretch}{2em}
\multlinegap0pt
\usepackage{amssymb}
\usepackage{enumerate}
\usepackage[shortlabels]{enumitem}
\newtheorem{lemma}{Lemma}
\newcommand{\Nat}{\mathbb N}
\newcommand{\Rea}{\mathbb R}
\title{Note on almost isometric ideals and local retracts in Banach and metric spaces}
\author{Leandro Candido, Marek C\'uth, Ond\v{r}ej Smetana}
\date{Source: arXiv:2312.07749v1 (CC BY 4.0)}
\begin{document}
\maketitle

\noindent\emph{Source and licence.} Note on almost isometric ideals and local retracts in Banach and metric spaces. Version arXiv:2312.07749v1 (CC BY 4.0), distributed by arXiv under the Creative Commons Attribution 4.0 International licence, http://creativecommons.org/licenses/by/4.0/, as recorded in the arXiv licence metadata for this exact version and shown on the abstract page https://arxiv.org/abs/2312.07749v1. Full licence notice: \texttt{licenses/arxiv-2312.07749-CC-BY-4.0.txt}.\par\smallskip

\noindent\textit{Editorial scope.} Counted unit: the article's lemma lem:calculations (source0.tex lines 352--371). The article's complete original proof of it is delivered with it (source0.tex lines 372--430, one \texttt{proof} environment of 59 lines). No mathematics is written, repaired or re-derived: the statement and the proof are the article's own lines, in the article's own notation, and no retained line is substituted or re-flowed.  No further source-local context is needed: every cross-reference of every retained line resolves inside the two delivered spans.  Nothing was omitted from any retained span, so the package carries no omission record.  No retained line contains a citation, so the wrapper carries no bibliography and no bibliography entry.  No retained line contains a figure environment, an image-inclusion command or drawing code, so no image is delivered and the package declares no asset dependency.  Editorial formatting only, in addition: an AMS \texttt{amsart} preamble that restates the article's own declarations and macros used by the retained lines, namely the environments \texttt{lemma} on separate counters, together with the standard packages those lines need.  The retained environments print in this document as Lemma 1 in order of appearance, whereas the article prints the same items with the numbering of its own full document.  The complete native source of the article as distributed in the arXiv e-print for this exact version is bundled unchanged as \texttt{sources/151412/source0.tex}.
\begin{lemma}\label{lem:calculations}
Let $X$ be a Banach space, $U,V\subset X$ finite-dimensional subspaces with $U\cap V = \{0\}$. Then for any $m\in\Nat$ there is a function $\zeta:[0,1)\to[0,1)$ continuous at zero with $\zeta(0)=0$ such that whenever $\varepsilon\in (0,1]$,
\begin{itemize}
    \item $n\in\Nat$ is such that $n\geq \tfrac{1}{\varepsilon}$,
    \item $N\subset V$ is an $\varepsilon$-net of $nB_V$,
    \item $A\subset \ell_1^{\dim U}$ is an $\varepsilon$-net of $S_{\ell_1^{\dim U}}$
    \item $\{e_1,\ldots,e_{\dim U}\}\subset mB_U$ is a basis of $U$,
    \item $\{f_1,\ldots,f_{\dim U}\}\subset mB_X$ are such that for every $x\in N$ and $a\in A$ we have
        \[
            \Big|\|x+\sum_{i=1}^{\dim U}a_i e_i\| - \|x+\sum_{i=1}^{\dim U}a_i f_i\|\Big| \leq \varepsilon,
        \]
\end{itemize}
then the mapping $T:U\oplus V\to X$ defined by 
\[
T(x+\sum_{i=1}^{\dim U}a_i e_i):=x+\sum_{i=1}^{\dim U}a_i f_i,\quad x\in V,\; a\in \Rea^{\dim U}
\]
is $(1+\zeta(\varepsilon))$-isomorphism.

Moreover, if $x^*\in X^*$ is such that $|x^*(e_i-f_i)|<\varepsilon$, then $\|(x^*\circ T - x^*)|_{U\oplus V}\| < \zeta(\varepsilon)\|x^*\|$.
\end{lemma}

\begin{proof}Let us denote by $P_U$ the projection $P_U:U\oplus V\to U$. Further, since all the norms on a finite-dimensional space are equivalent, pick $C>0$ such that $\|\sum_{i=1}^{\dim U} a_ie_i\|\geq C\|a\|_1$ for any $a\in\Rea^{\dim U}$.

Pick any $a\in S_{\ell_1^{\dim U}}$ and $x\in nB_V$. Let $b\in A$ and $y\in N$ be such that $\|a-b\|_1\leq \varepsilon$ and $\|x-y\|\leq \varepsilon$. Then we have
	\begin{equation} \label{key_lemma_aux_calc_3}
		\begin{split} 
			\Bigg| \Big\| x & + \sum_{i=1}^{\dim U} a_i e_i \Big\| - \Big\| x + \sum_{i=1}^{\dim U} a_i f_i \Big\| \Bigg| \\
			& \leq \| x - y \| + \Big\| \sum_{i=1}^{\dim U} a_i e_i - \sum_{i=1}^{\dim U} b_i e_i \Big\| + \|x-y\|  + \Big\| \sum_{i=1}^{\dim U} a_i f_i - \sum_{i=1}^{\dim U} b_i f_i \Big\| \\
			& \quad + \Bigg| \Big\| y + \sum_{i=1}^{\dim U} b_i e_i \Big\| - \Big\| y + \sum_{i=1}^{\dim U} b_i f_i \Big\| \Bigg| \\
			& < \varepsilon + \varepsilon m + \varepsilon + \varepsilon m + \varepsilon = \varepsilon(3 + 2m).
		\end{split}
	\end{equation}
Moreover, for any $z\in V$ and $c\in \Rea^{\dim U}$ we have
\[
    \Big\|z + \sum_{i=1}^{\dim U}c_i e_i\Big\| \geq \frac{1}{\|P_U\|}\Big\|\sum_{i=1}^{\dim U}c_i e_i\Big\|\geq \frac{C\|c\|_1}{\|P_U\|},
\]
so from \eqref{key_lemma_aux_calc_3} we obtain
\[
    \Bigg| \Big\| x + \sum_{i=1}^{\dim U} a_i e_i \Big\| - \Big\| x + \sum_{i=1}^{\dim U} a_i f_i \Big\| \Bigg|\leq \frac{\varepsilon(3+2m)\|P_U\|}{C}\Big\| x + \sum_{i=1}^{\dim U} a_i e_i \Big\|,
\]
which implies that for $f(\varepsilon):=\min\{\frac{\varepsilon(3+2m)\|P_U\|}{C},\tfrac{1}{2}\}$ we have
\begin{equation}
    (1-f(\varepsilon))\Big\| x + \sum_{i=1}^{\dim U} a_i e_i \Big\|\leq \Big\| x + \sum_{i=1}^{\dim U} a_i f_i \Big\|\leq (1+f(\varepsilon))\Big\| x + \sum_{i=1}^{\dim U} a_i e_i \Big\|
\end{equation}
and for $h(\varepsilon) = \tfrac{f(\varepsilon)}{1-f(\varepsilon)}$ we have $(1+h(\varepsilon))^{-1} = (1-f(\varepsilon))$ and therefore
\begin{equation}\label{key_lemma_final_ineq}
    \frac{1}{1+h(\varepsilon)}\Big\| x + \sum_{i=1}^{\dim U} a_i e_i \Big\|\leq \Big\| x + \sum_{i=1}^{\dim U} a_i f_i \Big\|\leq (1+h(\varepsilon))\Big\| x + \sum_{i=1}^{\dim U} a_i e_i \Big\|
\end{equation}
whenever $x\in nB_V$ and $a\in S_{\ell_1^{\dim U}}$.

    Pick arbitrary $a \in S_{\ell_1^{\dim U}}$ and $x \in X \setminus mB_U$. We observe
	\[
	\|x\| - m \leq \| x \| - \left\| \sum_{i=1}^{\dim U} a_i e_i \right\| \leq \left\| x + \sum_{i=1}^{\dim U} a_i e_i \right\| \leq \|x\| + m.
	\]
	The same argument yields
	\[
	\|x\| -m \leq \left\| x + \sum_{i=1}^{\dim U} a_i f_i \right\| \leq \|x\| + m.
	\]
	This means
	\begin{equation} \label{key_lemma_two_ineq}
		\frac{\left\| x + \sum_{i=1}^{\dim U} a_i f_i \right\|}{\left\| x + \sum_{i=1}^{\dim U} a_i e_i \right\|} \leq \frac{\|x\| + m}{\|x\| - m} \text{ and } \frac{\left\| x + \sum_{i=1}^{\dim U} a_i e_i \right\|}{\left\| x + \sum_{i=1}^{\dim U} a_i f_i \right\|} \leq \frac{\|x\| + m}{\|x\| - m}.
	\end{equation}
	We consider the real function $g: z \mapsto \frac{z+m}{z-m} - 1$. Because $\| x \| > n\geq \tfrac{1}{\varepsilon}$ and the function $g$ is decreasing to $0$, we have $\frac{\|x\| + m}{\|x\| - m} = g(\|x\|) +1<g(\tfrac{1}{\varepsilon}) + 1$. From this and \eqref{key_lemma_two_ineq} we obtain	\begin{equation}\label{eq:keyLemmaOdhadVetsiNezGJedna}
	\frac{1}{1 + g(\tfrac{1}{\varepsilon})} \Big\| x + \sum_{i=1}^{\dim U} a_i e_i \Big\|\leq \Big\| x + \sum_{i=1}^{\dim U} a_i f_i \Big\| \leq (1 + g(\tfrac{1}{\varepsilon})) \Big\| x + \sum_{i=1}^{\dim U} a_i e_i \Big\|.
	\end{equation}

 Thus, for $\zeta(\varepsilon):=2\max\{h(\varepsilon),g(\tfrac{1}{\varepsilon}),\frac{\varepsilon}{C} \|P_U\|\}$ from \eqref{key_lemma_final_ineq} and \eqref{eq:keyLemmaOdhadVetsiNezGJedna} we obtain that 
 \[
    \frac{1}{1 + \zeta(\varepsilon)} \Big\| x + \sum_{i=1}^{\dim U} a_i e_i \Big\| < \Big\| T\Big(x + \sum_{i=1}^{\dim U} a_i e_i\Big) \Big\| < (1 + \zeta(\varepsilon)) \Big\| x + \sum_{i=1}^{\dim U} a_i e_i \Big\|
 \]
 for every $a\in S_{\ell_1^{\dim U}}$ and every $x\in V$, which easily using a homogeneous argument implies that $T$ is $(1+\zeta(\varepsilon))$-isomorphism.

 In order to prove the ``Moreover'' part, pick $x^*\in X^*$ with $|x^*(e_i-f_i)|<\varepsilon$ and $y\in B_{U\oplus V}$. Let us denote $y = x + \sum_{i=1}^{\dim U} a_i e_i$. Then we have
	\begin{align*}
		\left| x^*(Ty) - x^*(y) \right| &= \left| x^*\left(x + \sum_{i=1}^{\dim U} a_i f_i\right) - x^*\left(x + \sum_{i=1}^{\dim U} a_i e_i\right) \right| \\
		&= \left| \sum_{i=1}^{\dim U} a_i x^*(e_i - f_i) \right| < \|x^*\|\sum_{i=1}^{\dim U} |a_i| \varepsilon = \|x^*\|\varepsilon \|a\|_{\ell_1^{\dim U}} \\
		&\leq \frac{\varepsilon\|x^*\|}{C} \left\| \sum_{i=1}^{\dim U} a_i e_i \right\| = \frac{\varepsilon\|x^*\|}{C} \| P_U(y) \| \leq \frac{\varepsilon\|x^*\|}{C} \|P_U\| \|y\|.
	\end{align*}
	Thus, we have that $\|(x^*\circ T - x^*)|_{U\oplus V}\| \leq \frac{\varepsilon}{C} \|P_U\|\|x^*\| < \zeta(\varepsilon)\|x^*\|$.
\end{proof}

\end{document}
